\section{现代 VLA 延伸：从 2023 Recipe 到统一模型}

本节是 2026 年的工程解读，不属于 2023 年课堂事实。它不使用后来系统反向证明本讲正确，而是分析 Ted Xiao 的 recipe 在统一 vision-language-action model 中仍然保留哪些约束。

\subsection{统一模型仍需要模块化验收}

即使把视觉、语言和动作放进一个模型，训练与部署仍应分别审计 perception、planning、skill execution、feedback 与 data coverage。端到端 loss 可以共享表示，却会让错误来源更难定位；因此系统层 evaluation 仍需保留 SayCan 式可执行性、Inner Monologue 式闭环和 DIAL 式数据语义审计。

\begin{importantbox}{统一参数不等于统一责任}
一个 checkpoint 可以同时承担多种功能，但工程契约仍要回答：哪类 token 能控制真实动作，谁验证动作安全，谁维护长期状态，谁判断任务完成，哪些输入只是数据而不是指令。参数共享不能替代权限与验收边界。
\end{importantbox}

\subsection{Data Mixture 需要显式 Resource Accounting}

现代多机器人数据混合应记录每个 embodiment、任务、环境与语言标签的比例。若一个高频简单任务占据大部分 token，训练 loss 下降可能掩盖稀有关键技能退化；若不同 action spaces 没有规范化，模型会把同一 token 学成冲突控制。RT-1 的 task-diversity 结果仍提醒我们：token count 不是覆盖率。

\subsection{Memory 与 Latency 是共同约束}

更长 context 可以保留更多视觉历史，却增加前向成本和控制延迟。具身系统通常需要把近期高频状态、长期任务摘要和可检索 episodic memory 分层，而不是把所有帧永久放入 attention window。评价时应同时报告 success、latency、control frequency、memory size 与 recovery behavior。

\begin{warningbox}{现代化解读的边界}
课堂系统拥有封闭厨房、有限动作和专用 detector；开放世界 VLA 面临更多传感器、人员安全和不可恢复错误。迁移的是“接口与证据纪律”，不是把课堂成功率直接外推到家庭或工业部署。
\end{warningbox}

\subsection{本章小结}

统一 VLA 可以吸收 2023 recipe 的多种模块，但仍需要独立可执行性、反馈、数据覆盖和实时预算验收。Foundation model 的统一表示越强，系统越应透明报告什么仍由专用 data、controller 或 safety layer 提供。

